% GENERATED FILE. Edit canonical lessons, then run npm run generate:books.
% canonical-source-hash: fnv1a64:e88908db8b506bb6
% canonical-lessons: MW-W02-aa-independent, MW-W02-bha, MW-C02-aabhaar, MW-C02-practice

\chapter{Aabhaar: Formal Thanks, Two New Signs}
\label{ch:mw-aabhaar}
\hlchaptermodality{2}

\begin{chapteropening}
\textbf{By the end of this chapter:} \emph{I can recognise, hear, say, copy, and independently write \mw{आभार} as a short formal gratitude response.}

A source-attested formal gratitude word built from two new signs and two signs already taught, followed by independent four-skill retrieval.
\end{chapteropening}

\section[ā]{\mw{आ} — long \emph{ā} at the beginning}

\label{lesson:MW-W02-aa-independent}



\begin{quote}
With the page covered, write \textbf{\mw{राम}} and give the greeting reply you learnt in Chapter 1. Uncover. Point to the vowel mark in \textbf{\mw{रा}}: you know \textbf{\mw{ा}} when long \emph{ā} follows a consonant.
\end{quote}

\subsection*{Script — one new sign}
When long \emph{ā} begins a word, it stands by itself:

\begin{quote}
\textbf{\mw{आ}} — \emph{ā}
\end{quote}

Notice the upright line at the right. Say one steady \emph{ā} while your finger follows the sign.

\subsection*{Writing — observe, trace, copy}
1. look at \textbf{\mw{आ}} for three seconds 2. finger-trace it once 3. copy it twice 4. compare only the upright line and the top bar

Do not add a word yet. This session belongs to one sign.

\begin{tcolorbox}[breakable,colback=teal!4,colframe=teal!35!black,title={Before you move on}]
Cover the models. Write \textbf{\mw{आ}} once, then uncover and compare.
\end{tcolorbox}
\section[bha]{\mw{भ} — \emph{bha}, with a small breath}

\label{lesson:MW-W02-bha}



\begin{quote}
Write \textbf{\mw{स}}, then \textbf{\mw{आ}}, from memory. Check both before moving on.
\end{quote}

\subsection*{Script — one new sign and sound}
\begin{quote}
\textbf{\mw{भ}} — \emph{bha}
\end{quote}

Begin like English \emph{b}, then let a small puff of air follow. Hold a hand near your mouth: \emph{bha}. The breath is part of the sound.

\subsection*{Writing — observe, trace, copy}
1. look at \textbf{\mw{भ}} 2. finger-trace its lower shape, then the upright line 3. copy \textbf{\mw{भ}} twice 4. say \emph{bha} after each copy

Stop there. The next lesson will join this sign to signs already in your hand.

\begin{tcolorbox}[breakable,colback=teal!4,colframe=teal!35!black,title={Before you move on}]
Cover the model for five seconds. Write \textbf{\mw{भ}}, uncover, and compare.
\end{tcolorbox}
\section[ābhār]{\mw{आभार} — formal thanks, built from known signs}

\label{lesson:MW-C02-aabhaar}



\begin{quote}
Write \textbf{\mw{सा}} from memory. Then read these four pieces: \textbf{\mw{आ} · \mw{भ} · \mw{ा} · \mw{र}}.
\end{quote}

\subsection*{Script — build the word}
Join only signs you already know:

\begin{quote}
\textbf{\mw{आ}} + \textbf{\mw{भ}} + \textbf{\mw{ा}} + \textbf{\mw{र}} = \textbf{\mw{आभार}}
\end{quote}

>

\begin{quote}
\emph{ābhār} — formal thanks; gratitude
\end{quote}

Rajasthan Sahitya Akademi's \emph{Jāgtī Jot} uses \textbf{\mw{आभार}} when expressing gratitude in Rajasthani prose. It is a formal, Sanskrit-derived word. This lesson keeps that register visible rather than pretending every kind of thanks sounds the same.

\subsection*{Writing — from guided to delayed copy}
1. copy \textbf{\mw{आभार}} once while looking 2. point to \textbf{\mw{भा}} and say \emph{bhā} 3. hide the word for ten seconds 4. write \textbf{\mw{आभार}} from memory 5. compare one sign at a time

If one sign stalls, practise only that sign before trying the word again.

\begin{tcolorbox}[breakable,colback=teal!4,colframe=teal!35!black,title={Before you move on}]
Read \textbf{\mw{आभार}}, say its meaning, then write it once with the model covered.

Source: \href{https://rsad.artandculture.rajasthan.gov.in/content/dam/doitassets/art-and-culture/Rajasthani-Bhasha-Sahitya-Avm-Sanskriti-Academy-Bikaner/pdf/JJ\_All\_Pdf/JJ\_33\_01\_01\_April\_to\_April\_2004.pdf}{Rajasthan Sahitya Akademi, \emph{Jāgtī Jot}, April 2004}.
\end{tcolorbox}
\section[Practice]{Practice — receive a kindness, offer formal thanks}

\label{lesson:MW-C02-practice}



\begin{quote}
Give the respectful Chapter 1 greeting from memory. Then, with the new word covered, write \textbf{\mw{आ}}, then \textbf{\mw{भ}}. Uncover and check.
\end{quote}

\subsection*{The exchange — one courtesy move}
Someone helps you or gives you something. Offer formal thanks:

\begin{quote}
\textbf{\mw{आभार।}} — \emph{Ābhār.} — Thank you. / My gratitude.
\end{quote}

Use it here as a complete short response. A later chapter can add a more casual Marwadi gratitude form; do not make this one word carry every register.

\subsection*{Your turn — four tiny skills}
1. \textbf{Listen:} ask someone or a screen reader to say \emph{ābhār}; identify it as formal gratitude. 2. \textbf{Speak:} imagine receiving help and say \textbf{\mw{आभार}} without reading. 3. \textbf{Read:} choose the thanks word from \textbf{\mw{राम} · \mw{आभार} · \mw{सा}}. 4. \textbf{Write:} close the book for fifteen seconds, then write \textbf{\mw{आभार}} from memory.

The written response counts only with the model hidden.

\begin{tcolorbox}[breakable,colback=teal!4,colframe=teal!35!black,title={Before you move on}]
If \textbf{\mw{आ}} or \textbf{\mw{भ}} stalled, return to that sign's two-minute practice before repeating the whole word.
\end{tcolorbox}
